\section{Nuclear Spectroscopic Measurements}

\begin{example}{2}{6-1: Radii in the Bohr Model}
\allbold{The Bohr radius, \(a_0\) is the radius of an electron in the ground state of a Hydrogen atom \[a_0=\frac{4\pi \epsilon_0 \hbar^2}{m_e e^2}\] The corresponding radius for the first Bohr orbit for an element with higher \(Z\) is \[r_Z = \frac{a_0}{Z}\] a.) What would be the radius of the first Bohr orbit for a muon bound to \ce{^{150}_{60}Nd}?}

For a muonic atom, replace \(m_e\rightarrow m_\mu\) to give \(a_0 = 2.55 \times 10^{-13}\). Then substitute as follows.
\[r_{60} = \frac{2.55 \times 10^{-13}}{60} = 4.25\,\text{fm}\]

\allbold{b.) Is this orbit, on average, `inside' the nucleus?}

Assuming a spherical nucleus, \(R=R_0A^{1/3} = 1.05\times 10^{-15} (150)^{1/3} = 5.57\,\text{fm}\). Since \(4.25 < 5.57\,\text{fm}\), yes this orbit is `inside' the nucleus.

\allbold{c.) Using the concept of the first and second Bohr orbits, can the energy scale in the following spectra of muonic atom transitions be simply explained?}
    
\end{example}

\begin{example}{1}{6-4: Precision of a Penning Trap}
\allbold{The precision on nuclear masses which can be obtained depends on the number of ions in the measurement, \(N\), and the mass resolving power, \(R\): \[\frac{\delta m}{m}\propto \frac{1}{R\sqrt{N}}\]
The mass resolving power for a Penning trap is typically \(10^5\) for an observation time of \(100\) ms. This means that the isotope under study must be sufficiently long-lived to allow a measurement to be conducted for \(100\) ms. Estimate the number of ions that would need to be trapped to achieve a mass precision, \(\delta m/m=10^{-7}\)}

Rearrange the above equation for \(N\).
\[N=\frac{1}{\left(R\dfrac{\delta m}{m}\right)^2} = \frac{1}{(10^5 \times 10^{-7})^2}=10^4\,\text{ions}\]
\end{example}

\begin{example}{3}{6-5: Lutetium-176}
\allbold{The isotope \ce{^{176}Lu} occurs naturally. Explain this in terms of expectations based on masses of \(A=176\) isobars and spins of states in \ce{^{176}Lu} and its neighbours.}
\end{example}

\begin{example}{1}{6-6: Isotopes of Uranium}
\allbold{The isotopes \ce{^{235}U} and \ce{^{238}U} both occur naturally. They possess relative natural abundances of \(0.7204\%\) and \(99.2742\%\), and half lives of \(7.04\times 10^{8}\) y and \(4.468\times 10^{9}\) y respectively. Assuming that these two isotopes were present in equal amounts when they were formed (in some coalescing binary neutron star pair) estimate how long ago that was. Compare this with the estimated age of the Earth, \(4.5\times 10^9\) y.}    
\end{example}


\begin{example}{1}{6-7: The Fraction of Plutonium in Uranium Ore}
\allbold{The longest lived isotope of plutonium is \ce{^{244}Pu} with a half life of \(8.13 \times 10^7\) y. Assuming that plutonium and uranium have similar mineral chemistry and were formed in equal amounts as in the scenario in 6-6, estimate the ratio \ce{^{244}Pu}/\ce{^{238}U} in uranium ore. Has this been investigated?}
\end{example}


Questions for David
What is the difference between an evaluation and a compilation of nuclear data

How do you choose the reference point for an isotope shift?
In the Bethe-Weizsacker formula (1-13) why is there no pairing term?
Removing the surface term - would there be much change for heavy nuclei? their B would increase, but comparatively not as much as small ones.

square vs curved brackets in the ensdf